\chapter{Blurred Vision}
\index{Blurred Vision}

\section*{Definition and Overview}
Blurred vision is a loss of sharpness or clarity of sight, making objects look out of focus. It is a common symptom with many possible causes, ranging from a simple need for glasses to serious conditions of the eye, brain, or body. Blurred vision can affect one or both eyes, come on suddenly or gradually, and may be constant or come and go. Because some causes are sight-threatening or are the first sign of a stroke or other emergency, blurred vision always deserves careful assessment rather than being dismissed as insignificant.

\section*{Epidemiology}
Blurred vision is one of the most common reasons people see an optometrist or doctor. Refractive error -- the need for glasses -- affects a very large proportion of the population and is the leading cause of correctable vision problems worldwide. Cataracts are the most common cause of reversible vision loss in older adults, while diabetic eye disease and glaucoma are important causes of permanent visual impairment. Sudden visual loss is less common but is an emergency, most often resulting from problems with the blood supply to the eye or the vision-processing areas of the brain.

\section*{Aetiology and Pathophysiology}
The causes of blurred vision are many and act at different points along the visual pathway:

\begin{itemize}
    \item \textbf{Refractive error:} the eye focuses light imperfectly because of its shape -- shortsightedness, longsightedness, and astigmatism
    \item \textbf{The front of the eye:} dry eye, a scratched or infected cornea, or a cataract clouding the lens
    \item \textbf{The retina:} damage to the light-detecting layer at the back of the eye, as in diabetes, age-related macular degeneration, or retinal detachment
    \item \textbf{The optic nerve:} glaucoma, or inflammation of the nerve itself (optic neuritis)
    \item \textbf{The brain:} migraine, stroke, or transient ischaemic attack affecting the vision-processing areas
    \item \textbf{Whole-body conditions:} high or low blood sugar, high blood pressure, pregnancy, and some medicines can cause temporary blurring
    \item \textbf{Eye strain and fatigue:} prolonged close work, poor lighting, or uncorrected refractive error
\end{itemize}

\section*{Clinical Features}
The pattern of symptoms gives important clues to the cause:

\begin{itemize}
    \item Difficulty focusing, seeing faces, or reading small print
    \item Onset may be gradual (as in cataract or needing glasses) or sudden (as in stroke or retinal detachment)
    \item One eye or both eyes affected
    \item May be accompanied by pain, redness, sensitivity to light, or headache
    \item Associated features such as flashing lights, floaters, or a curtain-like shadow suggest retinal problems
    \item General symptoms such as thirst, frequent urination, or headache may point to diabetes or migraine
\end{itemize}

\section*{Differential Diagnosis}
Blurred vision can be confused with, or caused by, many conditions:

\begin{itemize}
    \item A need for new glasses versus genuine eye disease
    \item Dry eye syndrome versus corneal infection
    \item Migraine with visual aura versus a transient ischaemic attack
    \item Retinal detachment versus vitreous detachment with floaters
    \item Cataract versus age-related macular degeneration in older adults
    \item Glaucoma versus other causes of gradual vision loss
    \item Diabetes-related retinal changes versus the effects of high blood pressure on the retina
\end{itemize}

\section*{When to Seek Medical Care}
See a doctor or optometrist if:

\begin{itemize}
    \item Vision is blurred for more than a few hours or is progressively worsening
    \item There is pain, redness, or discharge from the eye
    \item Blurring affects only one eye, or comes and goes
    \item There are flashes of light, new floaters, or a shadow across the vision
    \item You have diabetes, high blood pressure, or a family history of eye disease
    \item Blurring is associated with headache, nausea, or sensitivity to light
\end{itemize}

\textbf{Contact your local emergency services for sudden, severe visual loss or:}
\begin{itemize}
    \item Sudden blurred vision with weakness, numbness, or difficulty speaking (possible stroke)
    \item Sudden vision loss in one or both eyes
    \item Severe eye pain with vision loss
\end{itemize}

\section*{Diagnosis and Investigations}
Assessment of blurred vision involves:

\begin{itemize}
    \item \textbf{History:} onset, which eye is affected, associated symptoms, past medical and eye history, and current medicines
    \item \textbf{Visual acuity testing:} reading an eye chart to measure the sharpness of vision
    \item \textbf{Refraction:} determining whether glasses improve the vision
    \item \textbf{Slit-lamp examination:} a microscope to inspect the front of the eye for cataract, infection, or dryness
    \item \textbf{Eye pressure measurement:} to check for glaucoma
    \item \textbf{Retinal examination:} looking at the retina and optic nerve after pupil dilation, to detect diabetic changes, macular degeneration, or detachment
    \item \textbf{Imaging and blood tests:} occasionally needed, such as scans of the brain or eye, or blood sugar and blood pressure checks
\end{itemize}

\section*{Management}
Management depends on the cause:

\begin{itemize}
    \item \textbf{Glasses or contact lenses:} correct refractive error
    \item \textbf{Lubricating drops:} for dry eye
    \item \textbf{Cataract surgery:} removes the cloudy lens and restores sight
    \item \textbf{Diabetes and blood pressure control:} protects the retina from further damage
    \item \textbf{Medical or laser treatment:} for diabetic retinopathy, macular degeneration, and some forms of glaucoma
    \item \textbf{Emergency treatment:} retinal detachment and stroke require urgent specialist care
    \item \textbf{Referral:} to an ophthalmologist or optometrist for assessment and follow-up
    \item \textbf{Adjusting medicines:} if a drug is causing the blurring
\end{itemize}

\section*{Complications}
\begin{itemize}
    \item Permanent vision loss if conditions such as retinal detachment, glaucoma, or stroke are not treated promptly
    \item Falls and injuries, especially in older people with poor vision
    \item Reduced ability to drive, work, or manage daily tasks
    \item Social isolation and reduced quality of life from impaired sight
    \item Missed diagnosis of an underlying condition such as diabetes or high blood pressure
    \item Complications of treatment, including those of eye surgery
\end{itemize}

\section*{Prognosis and Outcomes}
The outlook depends on the cause. Many causes of blurred vision are easily corrected, and refractive error is fully corrected with glasses or contact lenses. Cataract surgery restores good vision in the great majority of people. With diabetic retinopathy and glaucoma, early treatment slows or stops further damage, while delayed treatment can lead to permanent sight loss. Sudden causes such as retinal detachment or stroke have the best outcomes when treatment begins quickly, so prompt assessment is important.

\section*{Prevention and Self-Care}
\begin{itemize}
    \item Have regular eye examinations, with the frequency guided by age, risk factors, and symptoms
    \item If you have diabetes, attend retinal screening and keep blood sugar and blood pressure well controlled
    \item Wear protective eyewear for work and sport
    \item Follow the 20-20-20 rule for screen use: every 20 minutes, look at something about 20 feet away for 20 seconds
    \item Stop smoking, and maintain a healthy diet and weight
    \item Use good lighting when reading and keep prescribed glasses current
    \item Never ignore sudden visual loss or pain in the eye -- seek urgent care
\end{itemize}

\section*{Sources and Further Reading}
The following sources provide reliable, up-to-date information on blurred vision and eye health:

\begin{itemize}
    \item NHS. \textit{NHS guidance on blurred vision and eye conditions.} https://www.nhs.uk/conditions/
    \item Mayo Clinic. \textit{Information on blurred vision and causes of vision loss.} https://www.mayoclinic.org/
    \item MedlinePlus. \textit{Blurred vision patient information.} https://medlineplus.gov/
    \item MSD Manual. \textit{Overview of vision problems and eye disorders.} https://www.msdmanuals.com/
    \item World Health Organization. \textit{Blindness and vision impairment fact sheets.} https://www.who.int/
    \item Centers for Disease Control and Prevention. \textit{Vision health and eye disease information.} https://www.cdc.gov/
\end{itemize}
